\cotefile{19}
\scanpage{81}
adjoint à droite $h$, on aura la suite \unsure{étendue}

\diagram{\begin{tikzcd}[ampersand replacement=\&,column sep=small,row sep=small]
\deleted{\unsure{\varphi_{!}}} \& \deleted{\unsure{\varphi^*|\varphi}} \arrow[d,equal] \& \deleted{\unsure{\varphi_*}} \arrow[d,equal] \& \& \& \\
\& f_! \& f^* \arrow[d,equal] \arrow[rr,no head,to path={-- ([yshift=2mm]\tikztostart.north) -| (\tikztotarget.north)}] \& f_* \arrow[d,equal] \& f^! \arrow[d,equal] \& f_? \arrow[d,equal]\\
\& \& g_! \& g^* \arrow[d,equal] \& g_* \arrow[d,equal] \& g^! \arrow[d,equal]\\
\& \& \& h_! \& h^* \& h_*
\end{tikzcd}}

etc. On trouve ainsi le résultat (de proche en proche)

\uline{Proposition} Soit $\mathcal U$ un univers, \deleted{\unsure{$\mathbb N$}} contenant
un ensemble infini. Pour \unsure{tt|tout} entier $n$,
il existe deux $\mathcal U$-topos $\unsure{E|\mathcal E}$, $F$
(de la forme $\hat C$, $\hat C'$ si on veut, avec
$C$, $C'$ des petites catégories) et une
suite de $n$-foncteurs \inserted{\deleted{\illegible{1}}, dont aucun n'est pl. fid.,} $f_1,f_2,f_3\ \ldots,f_n$
alternativement de $E$ dans $F$ et de $F$ dans $E$,
chacun étant l'adjoint à droite de celui
qui précède. \unsure{Celà|Ceci} si on veut, tous les $f_{\unsure{i|j}}$
y inclus $=$ $f_1$ et $f_n$, étant des morphismes
de topos -- et \unsure{$=$} des morphismes \unsure{stricts|direct}].

\uline{Question}. Peut-on trouver une suite infinie
dans les deux directions ? Il suffirait de
trouver $f$ et $g$, tels que $g$ soit à la fois adjoint
à gauche et \unsure{,|à} droite de $f$, i.e. tel que les

% archive sheet number: 80

\scanpage{82}
topos \deleted{\unsure{des}} \unsure{forment|formé|formés} des couples d'un
ens $A$ \deleted{\unsure{mu}} muni d'un \unsure{autom|endom} bijectif $w$,
et d'une rétraction de $A/w$ sur une
partie \inserted{$E$} de $A^w\subset A/w$. Le foncteur $f$
est le foncteur qui associe à ces
\unsure{fourbis} le \unsure{fourbis} ($E$ muni de $\mathrm{id}_E$, $E$, $\mathrm{id}_E$)\unsure{!|?}

% archive sheet number: 81

\scanpage{83}
Catégories fibrées

\margin{\unsure{ARCH 19|ARCH 1a}}

22 p

% archive sheet number: 82

\scanpage{84}
\uline{Sites fibrés}

\diagram{\begin{tikzcd}[ampersand replacement=\&,row sep=tiny]
\unsure{E|\mathcal E}\arrow[d]\\ B
\end{tikzcd}}

catégorie fibrée, avec les \deleted{foncteurs} \inserted{catégories} fibres
des sites, les foncteurs $f^*$ des morphismes
de sites [continus \unsure{suffira|suffit|suffisantes} pour certaines
choses). On met sur $E$ la top. la moins fine de
celles qui rend \unsure{continu|continus} les foncteurs $E_b\longrightarrow E$.

Les familles couvrantes sont décrites par VI 3.2.

Supposons $B,E\in\mathcal U$, \deleted{\unsure{tous et}} (pour simplifier) les
top. des $E_b$ \unsure{standart|A-endroit|Abondant}. Plongeons chaque
$E_b$ dans $\widetilde E_b$, on trouve une nouvelle catégorie
fibrée $\overline E$, : fibres les \inserted{topos} $\widetilde E_b$ ; un $f$-\deleted{\unsure{foncteur}} morphisme de
$F\in\unsure{\overline E_x|\widetilde E_x}$ dans $G\in\widetilde E_y$ est un morphisme

\margin{$\unsure{f}\quad x\longrightarrow y$}

\[
\text{dans }\widetilde E_x \qquad F\longrightarrow f^*(G)
\]

[en désignant aussi par $f^*$ l'ext.
de $f^*$ aux $\widetilde E_b$].

Déterminons le topos $\widetilde E$, \unsure{le|des|dans} \deleted{\unsure{uns|un}}
\unsure{dernier point|énoncé général} nous dit que \unsure{ses} \deleted{\unsure{sections}} \inserted{éléments} s'identifient

\[
\deleted{\widetilde E\simeq \underset{X\in\unsure{\mathrm{Ob}}\,E}{\underleftarrow{\lim}}
(E_{p(X)})/X = \overset{\text{\unsure{des fibres|des lieux}}}{\underset{X\in\mathrm{Ob}\, B}{\underleftarrow{\lim}}}E_x}
\]

\deleted{mais étant entendu} \unsure{,}
aux sections sur $E$ de la catégorie fibrée dont
la fibre, en $X\in\mathrm{Ob}\,E$, est la \deleted{\unsure{fibre}} \inserted{catégorie \uline{opposée} au} topos
associé au ``petit site'' défini par $X$ dans $E$, savoir

% archive sheet number: 83

\scanpage{85}
$(\unsure{E|\mathcal E}_{p(x)})_{/X}$ \deleted{$\unsure{\simeq E_x|\simeq\widetilde E_x}$},
\deleted{mais en \unsure{astreignant|assignant}} qui sont
cartésiennes \unsure{quand|qd} $X$ parcourt une fibre
(``petit site'') $E_x$. De façon précise

\[
\widetilde E^\circ \approxeq \underline{\Gamma}^{\mathrm{vert\text{-}cart}}
(\deleted{\unsure{B}}E,\ \deleted{\unsure{E}}X\mathrel{\unsure{\leadsto|\sim}}
{(\overline E)_{p(X)/X}}^\circ),
\]

où l'exposant ``vert-cart'' signifie qu'on
se borne aux sections qui sont cartésiennes
au dessus des flèches verticales. Utilisant
une formule d'associativité évidente, le
deuxième \unsure{ms|membre} s'interprète comme

\[
\underline{\Gamma}\bigl(B,\ x\leadsto
\underbrace{\underline{\Gamma}^{\mathrm{cart}}\bigl(E_x,\ X\leadsto
(\overline E_x)^\circ_{\!/X\unsure{\rbrace|3}}\bigr)}_{\displaystyle\widetilde E_x^\circ=(\overline E)_x^\circ}
\bigr)
\]

Donc

\[
\widetilde E\simeq\underline{\Gamma}(B,\overline E^\triangle)^\circ
\]

où $\overline E^\triangle$ désigne la catégorie fibrée déduite de
$\overline E$ en passant aux \deleted{\unsure{duales}} \inserted{opposées} des fibres de $\overline E$.

L'inclusion

\[
\underleftarrow{\mathrm{Lim}}\bigl(E,X\leadsto(\overline E)_{p(X)/X}\bigr)
\longrightarrow\overset{\widetilde E}{\overset{\Vert}{\underline{\Gamma}^{\mathrm{vert.cart}}}}
(E,\ldots)
\]

n'est alors autre \inserted{\unsure{mod signis|mod isom}} que l'inclusion

\[
\unsure{[|\underline{\Gamma}}\ \underleftarrow{\mathrm{Lim}}\,(B,\overline E^\triangle)^\circ
\simeq\underleftarrow{\mathrm{Lim}}\,(B,\overline E)^{\deleted{\unsure{\sharp|*|\circ|\triangle}}}
\subset\underline{\Gamma}(B,\overline E^\triangle)^\circ
\]

% archive sheet number: 84

\scanpage{87}
\margin{Catégories $\underline{\Gamma}$ et $\underleftarrow{\mathrm{Lim}}$}

\diagram{\begin{tikzcd}[ampersand replacement=\&,row sep=tiny]
\unsure{E|\mathcal E}\arrow[d]\\ B
\end{tikzcd}}

Catégorie fibrée

\diagram{\begin{tikzcd}[ampersand replacement=\&,row sep=large,column sep=large]
\&\underleftarrow{\mathrm{Lim}}\ E/B\arrow[r,hook,"i"]\arrow[d,"F_c^*"]\&\underline{\Gamma}(E/B)\arrow[d,"F^*"]\\
\underline{X'\in\mathrm{Ob}}\&\underleftarrow{\mathrm{Lim}}\ E'/B'\arrow[r,hook,no head,"i'"]\&\underline{\Gamma}(E'/B')
\end{tikzcd}}

\[
\begin{aligned}
\mathrm{Hom}(X',F_c^*(Y))&=\mathrm{Hom}(i'(X'),F^*(i(Y)))\\
&=\mathrm{Hom}(F_*\,i'(X'),i(Y))\\
&=\mathrm{Hom}(\underbrace{i_*F_*\,i'(X')}_{F_*^c(X')},Y)
\end{aligned}
\]

\textcircled{1} \uline{Si $\exists$ adjoint à g. $i_{\unsure{+|*}}$ de $i$}, et si $\exists$ $F_{\unsure{*}}$ adjoint
à g. de $F^*$ (p.ex. si $E$ est \uline{aussi} cofibrée
sur $B$, et les ``petites'' $\unsure{\underrightarrow{\lim}|\underline{\lim}|\underleftarrow{\lim}}$ existent dans les fibres)
alors $\exists$ adjoint à gauche $F_*^c$ de ${F}_c^*$.

\[
\begin{aligned}
\mathrm{Hom}(F_c^*(Y),X')&=\mathrm{Hom}(F^*(i(Y)),i'(X'))\\
&=\mathrm{Hom}(i(Y),F_+\,i'(X'))\\
&=\mathrm{Hom}(Y,\underbrace{i_+F_+\,i'}_{F_+^c}(X'))
\end{aligned}
\]

\textcircled{2} \uline{Si} $\exists$ \uline{adjoint à dr.} $i_+$ \uline{de} $i$, et si $\exists$ adj. à dr.
$F_+$ de $F^*$ (p.ex. \deleted{\unsure{les fibres de $E/B$}} si les petites $\underleftarrow{\lim}$ existent ds

% archive sheet number: 86

\scanpage{88}
les fibres de $\unsure{E|\mathcal E}/B$), alors $\exists$ adjoint à droite $F_+^c$ de $F_c^*$.

\textcircled{3} \uline{Prop}. Supposons que pour \unsure{tt|tout} $b\in\mathrm{Ob}\,B$, la limite
inductive $\underset{c\in\mathrm{Ob}\,{}^b/B}{\underline{\mathrm{Lim}}}\ f^*(X(c))$ dans $E_b$ existe,
et supposons que cette $\underline{\lim}$ commute
à tout \unsure{ch.gt|ch.g\textsuperscript{t}} de base $u^*$ ($a:b'\mathbin{\unsure{-|\to}}b$).
Alors le foncteur adjoint à gauche $\deleted{F}i_*^{\deleted{\unsure{c|\varepsilon}}}$
de \deleted{$\unsure{F'|E}$} $i:\underleftarrow{\mathrm{Lim}}\ E/B\longrightarrow\underline{\Gamma}(B/E)$ existe,
et est donné par

\[
\deleted{F}i_*^{\deleted{\unsure{c|\varepsilon}}}(X)(b)
=\underset{x\in\mathrm{Ob}\,{}^b/B}{\underrightarrow{\mathrm{Lim}}}\ f^*(X(x))
\]

[Le morphisme d'adjonction

\[
\deleted{F}i_*^{\deleted{\unsure{c|\varepsilon}}}\,\deleted{F}i_{\deleted{\unsure{*|c}}}^{\deleted{\unsure{*|\varepsilon}}}(X)\longrightarrow X
\]

étant le morphisme évident].

\inserted{Si de plus $E\longrightarrow X$ est cofibrant, on trouve des adjoints à gauche des $b_c^*:\underleftarrow{\mathrm{Lim}}\ E/B\mathbin{\unsure{-|\to}}E_b$}

% slanted
\margin{Énoncé dual pour la construction d'un adjoint \uline{à droite}, pour
$\boxed{(\underleftarrow{\mathrm{Lim}}\ E/B)^\circ\to\underline{\Gamma}(E^\triangle/B)^\circ}$
moyennant la comm. des $f^*$ aux petites $\underleftarrow{\lim}$. Alors les adj.}

\deleted{N.B. Ce critère ne fait pas double emploi avec
l'hyp. \unsure{ess|ass} de commutation des $u^*$
aux $\underrightarrow{\lim}$ n'étant pas dans la $=$
direction que celle que les $u^*$ aient des
adjts à g. $u_*$ -- qui implique \deleted{\unsure{nt}} que $u^*$ commute
aux $\underleftarrow{\lim}$ \ldots\ ].}

\margin{\uline{à droite} des $b_c^*:\underleftarrow{\mathrm{Lim}}\ E/B\longrightarrow E_b$ existent également et se construisent explicit\textsuperscript{t}, s'il en est ainsi des $f_*^{\unsure{s|+|*}}$}

% archive sheet number: 87

\scanpage{89}
\uline{Cor} Si de plus les $b_*$ existent \inserted{($b\in\mathrm{Ob}\,B$)}
[p.ex. si $p$ est aussi cofibrant, et les ``petites'' sommes
dans les $\unsure{E|\mathcal E}_b$ existent), alors \unsure{le|les} ``foncteur
\unsure{fibre|fibres}'' $\underleftarrow{\mathrm{Lim}}\ E/B\xrightarrow{b_c^*}E_b$ a un adj. à
gauche $b_*^c$.

\uline{Cor} Supposons que $\underline{\Gamma}$ ait une \deleted{petite} famille
de générateurs \inserted{stricts} $(X_\lambda)$, alors les $\unsure{i_*|i}(X_\lambda)$
forment une famille de générateurs \inserted{stricts} de
$\underleftarrow{\mathrm{Lim}}$. En particulier, si $p$ cofibrant, et si
les $E_b$ ont des petites sommes et de
petites familles de générateurs stricts, (de sorte
qu'il en \unsure{soi.|soit|est} de $=$ dans $\underline{\Gamma}$), alors les $=\mathrm{s}$
conditions sont remplies pour $\underleftarrow{\mathrm{Lim}}$, et on trouve
un syst. de générateurs \inserted{stricts} en prenant pour \unsure{tt|tout}
$b\in\mathrm{Ob}\,B$, un ens $\mathrm E_b$ de gén. stricts de $\mathrm E_b$,
et prenant $\displaystyle\bigcup_{b\in\mathrm{Ob}\,B}b_*^c(\mathrm E_b)$.

% archive sheet number: 88

\scanpage{90}
[ Si \inserted{petites} $\underrightarrow{\lim}$ existent \inserted{ds $\unsure{E|\mathcal E}_b$,} et les $f^*$ y commutent :

\[
L\xrightarrow{i}\underline{\Gamma}(E/B)\qquad\text{a un adjoint à gauche }\inserted{i_*}
\]

Si petites $\underleftarrow{\lim}$ existent ds $E_b$, et les $f^*$ y commutent,

\[
L^\circ\xrightarrow{j^\circ}\underline{\Gamma}(E^{\triangle}/B)
\]

a un adjoint à gauche, donc

\[
L\mathrel{\overset{j}{\unsure{\longrightarrow|\rule{1.5em}{.35pt}}}}\underline{\Gamma}(E^\triangle/B)^\circ
\]

a un adjoint à droite $j_+$

Si on cherche des générateurs de $L$, il y a
lieu d'utiliser $i_*$, car \deleted{pour les} si on a des
gén $X_\lambda$ de $\underline{\Gamma}(E/B)$, comme

\[
\mathrm{Hom}_{\underline{\Gamma}}(X_\lambda,i(Y))\simeq
\mathrm{Hom}_L(i_*(X_\lambda),Y),
\]

on voit que la famille des $i_*(X_\lambda)$ forme un
syst. de gén. de $L$. Or, pour avoir un \inserted{(\uline{petit})}
syst. de gén de $\underline{\Gamma}$, le cor. 2 \inserted{de Deligne} appliqué
à $E^\circ\mathbin{\unsure{-|\to}}B^\circ$ nous dit qu'il suffit que
$p$ co\uline{fibrant} et que les petites sommes
existent ds les $E_b$, et que dans ceux-ci
il y ait des \deleted{gén.} petits ens. gén.

\[
\left|\ \begin{aligned}
&\text{Petites }\underrightarrow{\mathrm{Lim}}\ \text{existent ds les }E_b,\ \text{les }f^*\text{ y commutent}\\
&\text{Petit ens de générateurs dans les }E_b
\end{aligned}\right.
\]

% archive sheet number: 89

\scanpage{91}
Si on cherche des cogénérateurs de $L$
i.e. des générateurs de $L^\circ$, il y a lieu
d'utiliser $j^\circ$ associé à ${\varphi}^\triangle$, et on trouve
comme condition suff. que $j_*^\circ$ existe et qu'il
y ait petit syst. de gén. de $\underline{\Gamma}({\unsure{E|\mathcal E}}^\triangle/B)$,
p.ex. que \uline{les petites} $\underleftarrow{\lim}$ \uline{existent dans}
\uline{les} $E_b$, que \uline{les foncteurs $f^*$} y \uline{commutent},
\uline{que les} $E_b^{\unsure{\cdot|{}}}$ \uline{aient des petits ens de cogénérateurs}]

% archive sheet number: 90

\scanpage{92}
\deleted{Si \unsure{dans}}

Si dans les catégories fibres, certains \inserted{types de} $\underrightarrow{\lim}$ ou $\underleftarrow{\lim}$
sont représentables, et si les $f^*$ y commutent,
alors les $=\mathrm{s}$ types de lim sont représentables
dans $\underleftarrow{\mathrm{Lim}}\ \unsure{E|\mathcal E}/B$, et les foncteurs fibres (plus
gén, les foncteurs chang\textsuperscript{ts} de base par \unsure{des|un}
$B'\mathbin{\unsure{-|\to}}B$) y commutent. Ainsi, si les $E_b$\unsure{/|{}}
sont préadditives (\unsure{r.|resp.} additives, abéliennes) \inserted{et si les $f^*$ ``respectent'' ces struct.} il en est
de $=$ de \deleted{$\unsure{\Gamma}$} $\underleftarrow{\mathrm{Lim}}\ E/B$, si $E_b$ est à \inserted{petites} $\underrightarrow{\lim}$ filtrantes
\unsure{ex.ctes|exactes}, et si les $\underleftarrow{\lim}$ finies y sont représentables,
\inserted{\uline{et} si les $f^*$ commutent à \unsure{tt ça|t.t. ça}, alors} kif-kif pour $\underleftarrow{\mathrm{Lim}}\ E/B$. De $=$, si les $E_b$
satisfont \inserted{\unsure{l'une}} les conditions a) b) c) \unsure{de l'exactitude|d'exactitude} d'
un $\mathcal U$-topos, \inserted{et les $f^*$ des morphismes de topos,} de $=$ $\underleftarrow{\mathrm{Lim}}\ E/B$. Pour pouvoir
conclure que si les $E_b$ sont des $\mathcal U$-topos
(resp. des $\mathcal U$-topos \unsure{abéliens|distribués}) \inserted{et les $f^*$ des morphismes de tels,} alors il en est de
$=$ de \deleted{$\unsure{\Gamma}$} $\underleftarrow{\mathrm{Lim}}\ E/B$, il reste à donner des
conditions assurant que cette catégorie est une
$\mathcal U$-catégorie, et admet une \inserted{petite} famille de gén.
Le premier pt sera assuré dès que $B$ est
``ess. petite''. Pour \unsure{le|les} deuxième, si on veut
appliquer le corollaire qui précède, on est

% archive sheet number: 91

\scanpage{93}
amené à supposer que \deleted{\unsure{le|les}} $p$ est \unsure{égal\textsuperscript{t}} cofibrant,
i.e. que les adjoints \uline{à gauche} des $f^*$ existent,
[qu'il n'y a pas lieu alors de noter $f_*$, mais
plutôt $f_!$], \inserted{si les $f^*$ sont associés à des \uline{morphismes} $\unsure{f_{!}|f:}$,}
On sait qu'il revient au $=$ de
dire que les $f^*$ commutent aux \inserted{petites} $\underleftarrow{\lim}$,
-- -- ce qui implique alors que ces dernières
se calculent fibre par fibre.

\deleted{\unsure{contre exemple} pour le cas où les $f^*$ ne commutent
pas aux $\underleftarrow{\lim}$ quelconques ?}

Dans le cas général, \unsure{où} $B$ est \deleted{\unsure{équivalente}} ``ess petite'',
il n'est cependant pas difficile de prouver
qu'il y a un petit ens de générateurs
-- à rédiger en détail ..... . Donc on
trouve bien des $\mathcal U$-topos resp $\mathcal U$-ab-topos.
Il y a donc des $\underleftarrow{\lim}$, mais leur
calcul en général, ne se faisant pas
terme à terme, est un peu mystérieux.

% archive sheet number: 92

\scanpage{95}
\margin{Dualité pour catégories fibrées et cofibrées}

Catégorie

\diagram{\begin{tikzcd}[ampersand replacement=\&] \unsure{E|\mathcal E}\arrow[d,"p"'] \\ B \end{tikzcd}}

sur une autre

1) La catégorie est \uline{fibrée} ou $p$ fibrant, si pour \deleted{\unsure{toute}} flèche $f:x\longrightarrow y$ de $B$, on peut définir $f^*:E_y\mathrel{\rule{1.4em}{.35pt}}E_x$, avec transitivité \unsure{essentielle|essent\textsuperscript{le}} :
\inserted{$\operatorname{Hom}_f(X,Y)\simeq\operatorname{Hom}_{\mathrm{id}_x}(X,f^*(Y))$}
On trouve un pseudo-foncteur

\[ F:B^\circ\longrightarrow(\mathrm{Cat}) \]

2) La catégorie est \uline{cofibrée}, ou $p$ cofibrant, si $p^\circ$ est fibrant, i.e. si pour toute $f:x\mathrel{\rule{1.4em}{.35pt}}y$ de $B$, existe un foncteur $f_*:E_x\mathrel{\rule{1.4em}{.35pt}}E_y$, avec transitivités essentielles : $\operatorname{Hom}_f(X,Y)\simeq\operatorname{Hom}(f_*(\unsure{x|X|X'}),Y)\rbrace$. Il donne naissance à un pseudo-foncteur

\[ F:B\longrightarrow(\mathrm{Cat}) \]

\uline{\unsure{Liens|Lien}} a) Si $p$ est fibrant, pour que (pour $f:x\mathrel{\rule{1.4em}{.35pt}}y$ donnée) $f_*$ existe, il f. et s. que $f^*$ ait un adjoint à g. Pour que $p$ soit cofibrant, il f. et s. que les $f^*$ aient tous des adjoints à g.

b) Si $p$ est cofibrant, pour que $f^*$ existe, il f. et s. que $f_*$ ait un adjoint à droite. Pour que $p$ soit fibrant, il f. et s. que tous les $f_*$ aient des adjoints à droite.

La catégorie \uline{opposée relative} d'une catégorie fibrée $p:E\mathrel{\rule{1.4em}{.35pt}}B$ se définit seulement à équivalence près, c'est une autre catégorie fibrée, associée au pseudo-foncteur

% archive sheet number: 94

\scanpage{96}
\[ F^{\mathrm{opp}}:B\longrightarrow(\mathrm{Cat}) \]

défini par

\[
\begin{aligned}
F^{\mathrm{opp}}(x)&=F(x)^\circ\\
F^{\mathrm{opp}}(f^{\unsure{\circ|'|{}}})&=F(f)^\circ
\end{aligned}
\]

\deleted{\unsure{et|lui}}

et formule \unsure{analogue} pour isom de transitivité.

Il n'est pas évident qu'il en existe une description intrinsèque à \uline{isom} près, en évitant \unsure{tt|tout} recours au pseudo-foncteur.

Pour que $(\deleted{\unsure{B}}{\unsure{E|\mathcal E}}/B)^{\mathrm{opp}}$ soit aussi cofibrée, il f. et s. que les foncteurs $f^*:E_y\mathrel{\rule{1.4em}{.35pt}}E_x$ aient des adjoints à \uline{droite} (condition opposée à celle pour que $E/B$ soit cofibrée).

\uline{À retenir} le $=$ pseudo-foncteur définit deux catégories (fibrées) bien distinctes sur $B$.

% archive sheet number: 95
